\section{Teorema Pemetaan Terbuka}\label{C0694}
Salah satu pertanyaan yang dapat (dan patut) diajukan tentang setiap kategori konkret ialah apakah setiap morfisme bijektif selalu merupakan
isomorfisme.  Dalam beberapa kategori jawabannya jelas \,\emph{ya} (misalnya, dalam $\cat{VEC}$).  Dalam kategori lain
jawabannya jelas \,\emph{tidak} (dalam $\cat{TOP}$, misalnya, tinjau pemetaan identitas dari bilangan
real bertopologi diskret ke bilangan real bertopologi biasa).  Dalam kategori-kategori yang muncul dalam
kajian analisis, aljabar sering menarik dengan penuh semangat ke satu arah, sementara topologi melawan dengan kuat dan menarik
ke arah yang lain.  Di sini pertanyaan tersebut dapat menjadi sangat menarik---dan jawabannya sama sekali tidak sepele.  Dalam kategori
$\cat{BAN_\infty}$ yang terdiri atas ruang Banach dan pemetaan linear kontinu, pertanyaan ini ternyata memiliki kedalaman yang menarik.  Dalam
kasus ini aljabar menang; jawabannya ya.  Hal ini ternyata merupakan akibat langsung dari \emph{teorema
pemetaan terbuka} yang termasyhur.  Ketika Anda menelaah pembuktian hasil berikut, perhatikan bahwa untuk pemetaan yang dibahas,
kelengkapan domain dan kelengkapan kodomain sama-sama esensial---dan dengan alasan yang sangat berbeda.

\begin{defn} Suatu pemetaan $f \colon X \sto Y$ antara ruang-ruang topologis disebut
 \index{terbuka!pemetaan}%
 \index{pemetaan!terbuka}%
\df{terbuka} jika memetakan himpunan terbuka ke himpunan terbuka; yaitu, $f$ merupakan pemetaan terbuka apabila $f^{\sto}(U)$ terbuka dalam $Y$
setiap kali $U$ terbuka dalam~$X$.
\end{defn}

\emph{Teorema pemetaan terbuka} menyatakan bahwa surjeksi linear kontinu antara ruang-ruang Banach selalu merupakan pemetaan terbuka.
Ada satu rincian teknis yang membuat pembuktian hasil ini agak rumit.  Misalkan $T$ suatu pemetaan linear kontinu
dari ruang Banach $B$ ke ruang linear bernorma~$V$, dan $R$ adalah citra oleh $T$ dari suatu bola terbuka di sekitar
titik asal dalam~$B$.  Kita perlu mengetahui bahwa jika tutupan $R$ memuat suatu bola terbuka di sekitar titik asal dalam $V$, maka
$R$ sendiri memuat bola yang sama. Akan memudahkan jika bagian informasi yang krusial ini kita pisahkan sebagai proposisi berikut.

\begin{notn} Dalam ruang linear bernorma $V$, kita menyatakan dengan
 \index{<unop@$(V)_r$ (bola terbuka berjari-jari $r$ di sekitar titik asal)}%
$(V)_r$ bola terbuka di sekitar titik asal dalam $V$ yang berjari-jari~$r$.
\end{notn}

\begin{prop}\label{prop_for_OMT} Misalkan $T\colon B \sto V$ suatu pemetaan linear kontinu dari ruang Banach ke ruang linear bernorma.
Jika $r$, $s > 0$ dan $(V)_s \subseteq \clo{T\bigl(\,(B)_r\,\bigr)}$, maka $(V)_s \subseteq T\bigl(\,(B)_r\,\bigr)$.
\end{prop}

\begin{proof}[\emph{Petunjuk pembuktian}] Pertama, buktikan bahwa tanpa mengurangi keumuman kita dapat mengasumsikan $r = s = 1$.  Dengan demikian,
hipotesis kita adalah $(V)_1 \subseteq \clo{T\bigl(\,(B)_1\,\bigr)}$.

Misalkan $z \in (V)_1$.  Pembuktian selesai jika kita dapat menunjukkan bahwa $z \in T\bigl(\,(B)_1\,\bigr)$.  Pilih $\delta > 0$ sedemikian
sehingga $\norm z < 1 - \delta$, dan misalkan $y = (1 - \delta)^{-1}z$.  Bangun suatu barisan $(v_k)_{k=0}^\infty$ dari vektor-vektor dalam~$V$
yang memenuhi
   \begin{enumerate}
     \item[(i)] $\norm{y - v_k} < \delta^k$ untuk $k = 0,1,2,\dots$; dan
     \item[(ii)] $v_k - v_{k-1} \in T\bigl(\,\delta^{k-1}(B)_1\,\bigr)$ untuk $k = 1,2,3,\dots$.
   \end{enumerate}
Untuk melakukannya, mulailah dengan $v_0 = \vc 0$. Gunakan hipotesis untuk menyimpulkan bahwa $y - v_0 \in \clo{T\bigl(\,(B)_1\,\bigr)}$ dan
karena itu terdapat vektor $w_0 \in T\bigl(\,(B)_1\,\bigr)$ sedemikian sehingga $\norm{(y - v_0) - w_0} < \delta$.  Misalkan
$v_1 = v_0 + w_0$, lalu periksa bahwa (i) dan (ii) berlaku untuk $k = 1$.

Kemudian cari $w_1 \in T\bigl(\,\delta (B)_1\,\bigr)$ sedemikian sehingga $\norm{(y - v_1) - w_1} < \delta^2$, dan misalkan $v_2 = v_1 + w_1$.
Periksa bahwa (i) dan (ii) berlaku untuk $k= 2$.  Lanjutkan secara induktif.

Untuk setiap $n \in \N$, pilih vektor $u_n \in \delta^{n-1}(B)_1$ sedemikian sehingga $T(u_n) = v_n - v_{n-1}$. Tunjukkan bahwa barisan
$(u_n)$ dapat dijumlahkan (gunakan~\ref{prop_abs_sum_implies_sum}), dan misalkan $x = \sum_{n=1}^\infty u_n$.  Selesaikan pembuktian dengan menunjukkan
bahwa $Tx = y$ dan bahwa $z \in T\bigl(\,(B)_1\,\bigr)$.   \ns
\end{proof}

\begin{thm}[Teorema pemetaan terbuka]\label{C069414} Setiap surjeksi linear terbatas antara
 \index{terbuka!pemetaan!teorema}%
 \index{morfisme bijektif!dapat diinverskan!dalam $\cat{BAN_\infty}$}%
ruang-ruang Banach merupakan pemetaan terbuka.
\end{thm}

\begin{proof}[\emph{Petunjuk pembuktian}] Misalkan $T\colon B \sto C$ suatu surjeksi linear terbatas antara ruang-ruang Banach.
Mulailah dengan membuktikan klaim berikut: yang perlu kita tunjukkan hanyalah bahwa $T\bigl((B)_1\bigr)$ memuat suatu bola terbuka di sekitar
titik asal dalam~$C$.  Untuk sementara, misalkan kita mengetahui bahwa bola semacam itu ada.  Namai bola itu~$W$.  Diberikan suatu subhimpunan terbuka $U$ dari $B$, kita
ingin menunjukkan bahwa citranya oleh $T$ terbuka dalam~$C$.  Untuk itu, ambil sebarang titik $y$ dalam $T(U)$ dan pilih $x$
dalam $U$ sedemikian sehingga $y = Tx$.  Karena $U - x$ merupakan lingkungan titik asal dalam $B$, kita dapat mencari bilangan $r > 0$ sedemikian sehingga
$r\,(B)_1 \subseteq U - x$.  Buktikan bahwa $T(U)$ terbuka dengan menunjukkan bahwa $y + rW \subseteq T(U)$.

Untuk membuktikan klaim di atas, misalkan $V = T\bigl((B)_1\bigr)$ dan perhatikan bahwa karena $T$ surjektif, $C = T(B) =
\bigcup_{k=1}^\infty k\clo V$.  Gunakan versi \emph{teorema kategori Baire} yang menyatakan bahwa jika suatu ruang metrik lengkap takkosong
merupakan gabungan suatu barisan himpunan tertutup, maka sekurang-kurangnya salah satu himpunan tersebut memiliki interior takkosong.
(Lihat, misalnya, Teorema 29.3.21 dalam~\cite{Erdman:2007}.)  Pilih $k \in \N$ sedemikian sehingga $\bigl(k\clo V\bigr)^o \neq \emptyset$.
Dengan demikian, terdapat $y \in \clo V$ dan $r > 0$ sedemikian sehingga $y + (C)_r \subseteq k\clo V$.  Dari kekonveksan
$k\clo V$, tidak sulit melihat bahwa $(C)_r \subseteq k\clo V$.  (Tuliskan sebarang titik $z$ dalam $(C)_r$ sebagai
$\frac12(y + z) + \frac12(-y + z)$.) Sekarang diperoleh $(C)_{r/k} \subseteq \clo V$.  Gunakan Proposisi~\ref{prop_for_OMT}.
\ns \end{proof}

\begin{cor}\label{C069417} Setiap bijeksi linear terbatas antara ruang-ruang Banach merupakan suatu
isomorfisme.
\end{cor}

\begin{cor} Setiap surjeksi linear terbatas antara ruang-ruang Banach
 \index{BAN@$\cat{BAN_\infty}$!hasil bagi dalam}%
 \index{hasil bagi!pemetaan!dalam $\cat{BAN_\infty}$}%
merupakan pemetaan hasil bagi dalam~$\cat{BAN_\infty}$.
\end{cor}

\begin{exam} Misalkan $\fml C_u$ himpunan fungsi kontinu bernilai real pada $[0,1]$ dengan
norma seragam, dan $\fml C_1$ himpunan fungsi yang sama dengan norma $\mathrm{L_1}$
(\,$\norm{f}_1 = \int_0^1\abs{f}\,d\lambda$\,). Maka pemetaan identitas $I\colon \fml C_u
\sto \fml C_1$ merupakan bijeksi linear kontinu, tetapi tidak terbuka. (Mengapa hal ini tidak
bertentangan dengan \emph{teorema pemetaan terbuka}?)
\end{exam}

\begin{exam} Misalkan $l$ himpunan barisan bilangan real yang pada akhirnya bernilai nol,
$V$ adalah $l$ yang dilengkapi norma $l_1$, dan $W$ adalah $l$ dengan norma seragam.  Pemetaan
identitas $I\colon V \sto W$ merupakan bijeksi linear terbatas, tetapi bukan suatu isomorfisme.
\end{exam}

\begin{exam} Pemetaan $T\colon \R^2 \sto \R\colon (x,y) \mapsto x$ menunjukkan bahwa meskipun
surjeksi linear terbatas antara ruang-ruang Banach memetakan himpunan terbuka ke himpunan terbuka, pemetaan tersebut tidak harus
memetakan himpunan tertutup ke himpunan tertutup.
\end{exam}

\begin{prop} Misalkan $V$ suatu ruang vektor, dan misalkan $\norm{\cdot}_1$ dan
$\norm{\cdot}_2$ adalah norma-norma pada~$V$ dengan topologi masing-masing $\ofml T_1$
dan~$\ofml T_2$.  Jika $V$ lengkap terhadap kedua norma dan jika $\ofml T_1
\supseteq \ofml T_2$, maka $\ofml T_1 = \ofml T_2$.
\end{prop}

\begin{exer} Apakah terdapat suatu barisan $(a_n)$ dari bilangan kompleks yang memenuhi
syarat berikut: suatu barisan $(x_n)$ dalam $\C$ dapat dijumlahkan secara mutlak jika dan hanya jika
$(a_nx_n)$ terbatas? \emph{Petunjuk.} Tinjau $ T\colon l_\infty \sto l_1\colon (x_n)
\mapsto (x_n/a_n)$.
\end{exer}

\begin{exam}\label{C069431} Misalkan $\fml C^2([a,b])$ keluarga semua fungsi bernilai real yang dua kali
dapat didiferensialkan secara kontinu pada interval $[a,b]$. Pandang keluarga ini sebagai ruang vektor
dengan cara biasa. Untuk setiap $f \in \fml C^2([a,b])$, definisikan
    \[ \norm{f} = \norm{f}_u + \norm{f'}_u + \norm{f''}_u\,. \]
Ini memang suatu norma, dan dengan norma ini $\fml C^2([a,b])$ merupakan ruang Banach.
\end{exam}

\begin{exam}\label{C069432} Misalkan $\fml C^2([a,b])$ ruang Banach yang diberikan dalam
Contoh~\ref{C069431}, dan misalkan $p_0$ dan $p_1$ merupakan anggota $\fml C([a,b])$.   Definisikan
  \[T\colon C^2([a,b]) \sto C([a,b])\colon
                      f \mapsto f'' + p_1f' + p_0f.\]
Maka $T$ merupakan pemetaan linear terbatas.
\end{exam}

\begin{exer} Tinjau suatu persamaan diferensial berbentuk
   \[ y'' + p_1\,y' + p_0\,y = q  \tag{$\ast$} \]
dengan $p_0$, $p_1$, dan $q$ fungsi-fungsi bernilai real kontinu (yang tetap) pada interval
$[a,b]$.  Nyatakan secara tepat, lalu buktikan, pernyataan bahwa \emph{solusi-solusi $(\ast)$
bergantung secara kontinu pada nilai awal}.  Anda boleh menggunakan tanpa bukti teorema standar berikut:
untuk setiap titik $c$ dalam $[a,b]$ dan setiap pasangan bilangan real $a_0$ dan $a_1$,
terdapat solusi tunggal untuk $(\ast)$ sedemikian sehingga $y(c) = a_0$ dan $y'(c) = a_1$.
\emph{Petunjuk.} Untuk $c \in [a,b]$ yang tetap, tinjau pemetaan
  \[ S\colon \fml C^2([a,b]) \sto \fml C([a,b]) \times \R^2 \colon
                                        f \mapsto \bigl(Tf,f(c),f'(c)\bigr) \]
dengan $\fml C^2([a,b])$ ruang Banach dari Contoh~\ref{C069431} dan $T$ pemetaan
linear terbatas dari Contoh~\ref{C069432}.
\end{exer}

\begin{defn} Suatu pemetaan linear terbatas $T\colon V \sto W$ antara ruang-ruang linear bernorma disebut
 \index{terbatas!jauh dari nol}%
\df{terbatas jauh dari nol} (atau
 \index{terbatas!dari bawah}%
\df{terbatas dari bawah}) jika terdapat bilangan $\delta > 0$ sedemikian sehingga $\norm{Tx} \ge
\delta\norm{x}$ untuk setiap $x \in V$. (Ingat bahwa kata ``terbatas'' memiliki satu arti ketika
menerangkan suatu pemetaan linear dan arti yang sangat berbeda ketika diterapkan pada fungsi umum;
demikian pula ungkapan ``terbatas jauh dari nol''. Lihat Definisi~\ref{def_bdd_away_zero}.)
\end{defn}

\begin{prop} Suatu pemetaan linear terbatas antara ruang-ruang Banach bersifat terbatas jauh dari nol jika dan
hanya jika pemetaan itu memiliki jangkauan tertutup dan kernel nol.
\end{prop}


















\section{Teorema Graf Tertutup}
Syarat perlu dan cukup yang sederhana agar suatu pemetaan linear $T\colon B \sto C$ antara ruang-ruang Banach
bersifat kontinu ialah grafnya tertutup.  Dalam satu arah, hal ini hampir jelas.  Kebalikannya
disebut \emph{teorema graf tertutup}.  Ketika kita mensyaratkan graf $T$
 \index{graf!tertutup}%
``tertutup'', maksud kita ialah tertutup sebagai subhimpunan $B \times C$ dengan topologi hasil kali.  Ingat (dari~\ref{exam_prod_top_norm})
bahwa topologi pada jumlah langsung $B \oplus C$ yang diinduksi oleh norma hasil kali sama dengan topologi hasil kali pada $B \times C$.

\begin{exer} Jika $T\colon B \sto C$ suatu pemetaan linear kontinu antara ruang-ruang Banach, maka grafnya tertutup.
\end{exer}

\begin{thm}[Teorema graf tertutup]\label{thm_CGT} Misalkan $T\colon B \sto C$ suatu pemetaan linear antara ruang-ruang Banach.
 \index{tertutup!teorema graf}%
Jika graf $T$ tertutup, maka $T$ kontinu.
\end{thm}

\begin{proof}[\emph{Petunjuk pembuktian}] Misalkan $G = \{(x,Tx)\colon x \in B\}$ graf~$T$. Terapkan
(Korolar~\ref{C069417} dari) \emph{teorema pemetaan terbuka} pada pemetaan
  \[ \pi\colon G \sto B\colon (x,Tx) \mapsto x\,. \]
\ns \end{proof}

Proposisi berikut memberikan syarat perlu dan cukup yang sangat sederhana agar graf suatu pemetaan linear antara ruang-ruang Banach
tertutup.

\begin{prop}\label{prop_nasc_closed_graph} Misalkan $T\colon B \sto C$ suatu pemetaan linear antara ruang-ruang Banach. Maka $T$
kontinu jika dan hanya jika syarat berikut dipenuhi:
   \[ \text{jika $x_n \sto 0$ dalam $B$ dan $Tx_n \sto c$ dalam $C$, maka $c = 0$.} \]
\end{prop}

\begin{prop} Misalkan $T \colon B \sto C$ suatu pemetaan linear antara ruang-ruang Banach sedemikian sehingga jika
$x_n \sto 0$ dalam $B$, maka $(g \circ T)(x_n) \sto 0$ untuk setiap $g \in C^*$.  Maka $T$
terbatas.
\end{prop}

\begin{prop} Misalkan $T \colon B \sto C$ suatu fungsi linear antara ruang-ruang Banach. Definisikan
$T^*f = f \circ T$ untuk setiap $f \in C^*$.  Jika $T^*$ memetakan $C^*$ ke dalam $B^*$, maka $T$
kontinu.
\end{prop}

\begin{defn}\label{C069624} Suatu pemetaan linear $T$ dari ruang vektor ke ruang itu sendiri disebut
 \index{idempoten}%
\df{idempoten} jika $T^2 = T$.
\end{defn}

\begin{prop}\label{prop_cont_idem_op} Suatu pemetaan linear idempoten dari ruang Banach ke ruang itu sendiri
kontinu jika dan hanya jika kernel dan jangkauannya tertutup.
\end{prop}

\begin{proof}[\emph{Petunjuk pembuktian}] Untuk bagian ``jika'', gunakan Proposisi~\ref{prop_nasc_closed_graph}.  \ns
\end{proof}

Proposisi berikut merupakan penerapan \emph{teorema graf tertutup} yang berguna pada operator ruang Hilbert.

\begin{prop} Misalkan $H$ suatu ruang Hilbert, serta $S$ dan $T$ fungsi dari $H$ ke dirinya sendiri sedemikian sehingga
    \[ \langle Sx,y \rangle = \langle x,Ty \rangle \]
untuk setiap $x$, $y \in H$.  Maka $S$ dan $T$ merupakan operator linear terbatas.
\end{prop}

\begin{proof}[\emph{Petunjuk pembuktian}] Untuk membuktikan linearitas pemetaan $S$, tinjau hasil kali dalam dari vektor-vektor
$S(\alpha x + \beta y) - \alpha Sx - \beta Sy$ dan $z$, dengan $x$, $y$, $z \in H$ dan $\alpha$, $\beta \in \K$.
Untuk membuktikan kekontinuan $S$, tunjukkan bahwa grafnya tertutup.  \ns
\end{proof}

Ingat bahwa dalam proposisi sebelumnya, operator $T$ adalah \emph{adjoin (ruang Hilbert)} dari $S$;
yaitu, $T = S^*$.

\begin{defn} Misalkan $a < b$ dalam~$\R$. Suatu fungsi $f\colon [a,b] \sto \R$ disebut
 \index{dapat didiferensialkan secara kontinu}%
 \index{diferensiabel!secara kontinu}%
dapat didiferensialkan secara kontinu jika fungsi itu dapat didiferensialkan pada (suatu himpunan terbuka yang memuat) $[a,b]$
dan turunannya $f\,'$ kontinu pada~$[a,b]$.  Himpunan semua fungsi bernilai real yang dapat
didiferensialkan secara kontinu pada $[a,b]$ dinotasikan dengan $\fml C^1(\,[a,b]\,)$.
\end{defn}

\begin{exam} Misalkan $D\colon \fml C^1(\,[0,1]\,) \sto C(\,[0,1]\,)\colon f \mapsto f\,'$, dengan
baik $C^1(\,[0,1]\,)$ maupun $C(\,[0,1]\,)$ dilengkapi norma seragam. Pemetaan
$D$ linear dan memiliki graf tertutup, tetapi tidak kontinu.  (Mengapa hal ini tidak bertentangan dengan
\emph{teorema graf tertutup}?  Apa yang dapat kita simpulkan tentang~$\fml C^1(\,[0,1]\,)$?)
\end{exam}

\begin{prop} Misalkan $S\colon A \sto B$ dan $T\colon B \sto C$ pemetaan linear antara
ruang-ruang Banach, serta $T$ terbatas dan injektif.  Maka $S$ terbatas jika dan hanya jika $TS$ terbatas.
\end{prop}











